\section{Preliminaries}

%
%
%
%

A \emph{bimatrix game} $\G = (M_{\mathrm{row}}, M_{\mathrm{col}})$ is
a game defined by two $N \times N$ matrices, $M_{\mathrm{row}}$
and $M_{\mathrm{col}}$, and two players: the \emph{row player} and the
\emph{column player}. We assume throughout that the game is normalized
so that both matrices have values in the interval $[0,1]$. The row and
column players choose \emph{strategies} $x$ and $y$ respectively,
where $x,y$ are nonnegative $N$-dimensional vectors satisfying
$\sum_ix_i=\sum_jy_j=1$.  A \emph{pure} strategy is one with support
$1$ (\ie, a vector with $1$ in one entry and $0$ in the rest). The
row (resp.\ column) player's \emph{payoff} is given by
$x^{\top}M_{\mathrm{row}}y$ (resp.\ $x^{\top}M_{\mathrm{col}}y$).

A \emph{Nash equilibrium} is a pair of strategies $(x,y)$ such that
neither player has any incentive to deviate to a different strategy,
assuming the other player does not deviate. Formally, in an
equilibrium, for all $i,j$ we have $e_i^{\top}M_{\mathrm{row}}y\leq
x^{\top}M_{\mathrm{row}}y$ and $x^{\top}M_{\mathrm{col}}e_j\leq
x^{\top}M_{\mathrm{col}}y$. An $\eps$-approximate Nash equilibrium, or
\emph{$\eps$-equilibrium} for short, is a pair of strategies $x,y$
where each player has incentive at most $\eps$ to deviate. That is,
for all $i,j$, 
\[
e_i^{\top}M_{\mathrm{row}}y\leq
x^{\top}M_{\mathrm{row}}y+\eps\qquad\text{and}\qquad
x^{\top}M_{\mathrm{col}}e_j\leq x^{\top}M_{\mathrm{col}}y+\eps\,.
\]

The \emph{value} of a pair of strategies, denoted $v_{\G}(x,y)$, is
the average payoff of the two players, \ie, 
\[
v_\G(x,y) = \frac{1}{2}(x^{\top}M_{\mathrm{row}}y +
x^{\top}M_{\mathrm{col}}y) = \sum_{i,j} x_i y_j
\frac{M_{\mathrm{row}}(i,j) + M_{\mathrm{col}}(i,j)}{2}\,.
\]

For a vector $x \in \Real^n$ and $S \subseteq [n]$, we write $x_S$ for
the projection of $x$ to the coordinates $S$.  We write $\|x\| =
\sum_{i = 1}^n |x_i|$ for the $\ell_1$ norm of $x$.  Thus, for a
strategy (in other words, a probability distribution) $x \in [0,1]^n$
we write $\|x_S\|$ for the probability that the player plays an
element of $S$.

Further, for a set $S \subseteq [n]$ of strategies, we use
$v_{\G|S}(x,y)$ to denote the value of $(x,y)$ \emph{conditioned} on
both players playing in $S$.  Formally,
\[
v_{\G|S}(x,y) = \E_{i \sim x, j \sim y}\left[
  \frac{M_{\mathrm{row}}(i,j)+M_{\mathrm{col}}(i,j)}{2}\;\middle|\;
  i,j \in S\right] = \frac{v_\G(x_S, y_S)}{\|x_S\| \cdot \|y_S\|}\,.
\]
(If $\|x_S\| = 0$ or $\|y_S\| = 0$, $v_{\G|S}(x,y)$ is undefined.)

Given an undirected graph $G=(V,E)$, and (not necessarily disjoint) vertex sets $S_1,S_2\subseteq V$, we will denote by $E(S_1,S_2)$ the set of ordered pairs $\{(i,j)\in S_1\times S_2\mid \{i,j\}\in E\text{ or }i=j\}$. We will refer to $d(S_1,S_2)=|E(S_1,S_2)|/(|S_1||S_2|)$ as the \emph{density} of the pair $(S_1,S_2)$.  Note that, letting $A$ be the adjacency matrix of $G$  (with $1$'s on the diagonal) and $\overline{x}, \overline{y} \in [0,1]^n$ be the uniform probability distributions over $S_1$ and $S_2$ (respectively), we have $d(S_1, S_2) = \overline{x}^{\top} A \overline{y}$.

We will write $\HC{n}{k}$ to denote the distribution in the input to the $k$-Hidden Clique problem. That is, we will write $G\sim\HC{n}{k}$ to mean that $G$ is sampled from $G(n,1/2)$ with a hidden clique of size $k$.

Throughout, w.\,h.\,p.~stands for ``with high probability,'' meaning $1-o(1)$.  Finally, we
shall make repeated use of the standard Chernoff-Hoeffding
bound~\cite{Hoeff63}.

\begin{lemma}[Chernoff-Hoeffding bound]
  \label{lemma:chernoff}
  Let $X_1, \ldots, X_m$ be independent random variables in $[0,1]$, and let
  $\mu = ({1}/{m}) \E[ \sum X_i]$.  Then, for all $\eps \ge 0$,
  $$
  \Pr\left[ \frac{1}{m} \sum X_i \ge \mu + \eps \right] \le e^{-2\eps^2 m}\,.
  $$
\end{lemma}


\section{Approximate equilibria with near-optimal value}
\label{sec:goodvalue}

\subsection{The reduction}\label{sec:reduction}

In this section we describe the general reduction that we use to prove
Theorems~1.2 and~1.4 and
describe its properties. This reduction also forms the basis for the reductions we use to prove Theorems~1.5, 1.6  and~1.9. It is based on the reduction of~\cite{HK11}.

As in~\cite{HK11} our soundness analysis proceeds by using the approximate equilibrium to find a dense
bipartite subgraph of $G$. The following lemma shows %
that this is sufficient to recover the hidden clique.

\begin{lemma}[{\cite[Lemma 5.3]{HK11}}]
  \label{lemma:reconstruct}
  There exist universal constants $c_1$ and $c_2$ such that the
  following holds.  Let $G\sim\HC{n}{C\ln n}$ %
%
 for some $C \ge c_1$.  Then,
  there is a polynomial time algorithm which, given a pair
  of vertex sets $S_1,S_2 \subseteq [n]$ of size $c_2 \ln n$ and
  density $d(S_1,S_2)\geq5/9$, reconstructs the hidden clique (with
  high probability over $G$).
\end{lemma}

The lemma is slightly different from Lemma 5.3 of~\cite{HK11}: there
we start with a bipartite subgraph of density $3/5$ instead of $5/9$
but this minor difference only changes the value of the constant $c_2$---the lemma holds for any constant density strictly larger than
${1}/{2}$.

Let us now describe the reduction.  It is controlled by three
parameters $\alpha, \beta, \gamma \in (0,1)$.  Setting these parameters
  appropriately gives the various hardness results.

\boxit{%
\begin{reduction}
  \label{red:reduction}
  Let $G = (V,E)$ be an $n$-vertex graph and $A$ its adjacency matrix
  (with $1$'s on the diagonal).  Then, for parameters $\alpha, \beta, \gamma \in (0,1)$, we define a (random) bimatrix game $\G := \G(G,
    \alpha, \beta, \gamma)$ as follows.

  Let $N = n^c$ where $c = (c_2+1) \ln (1/\beta)$ for the universal
  constant $c_2$ of \expref{Lemma}{lemma:reconstruct}.  Pick a random $N
  \times n$ matrix $B$ whose entries are i.\,i.\,d.\ $\{0,1\}$ variables
  with expectation $\beta$.  Then $\G = (M_{\mathrm{row}},
  M_{\mathrm{col}})$, where the payoff matrices are:
  \begin{align}\label{eq:simple-payoffs}
    M_{\mathrm{row}}=\left(
    \begin{array}{cc}
      \alpha A & 0\\
      B &\gamma J
    \end{array} 
    \right)
    & &
    M_{\mathrm{col}}=\left(
    \begin{array}{cc}
      \alpha A & B^{\top}\\
      0 &\gamma J
    \end{array} 
    \right),
  \end{align}
  and $J$ is the all-ones $N \times N$ matrix.
\end{reduction}
}

%
%
%

We conclude this section with an additional lemma 
which shows how to obtain a dense bipartite subgraph given an
approximate equilibrium of $\G$ with certain properties.  This lemma
(and its proof) is analogous to~\cite[Lemma 5.2]{HK11}, but as we need
it in larger generality we also give the proof.

\begin{lemma}
  \label{lem:dense-subgraph}
  Let $\G$ be as in \expref{Reduction}{red:reduction}.  Fix any $s \in
  [0,1]$, $t \in [0,1]$ and $\eps \in [0,1]$ such that $1 - 2t - 3
  \sqrt{s}/2 \ge \alpha + \eps$, and let $(x, y)$ be an
  $\eps$-approximate equilibrium of $\G$ with the following two
  properties:
  \begin{itemize}
  \item Both $\|x_{[n]}\| \ge 1-t$ and $\|y_{[n]}\|
    \ge 1-t$.
  \item The conditional value $v_{\G|[n]}(x,y) \ge (1-s)\alpha$.
  \end{itemize}
  Then, given $(x,y)$ as above, we can find, in polynomial time, vertex sets
  $S_1, S_2 \subseteq [n]$ each of size $c_2 \ln n$ such that their density is $d(S_1,S_2)\ge 5/9$.%
\end{lemma}

In the proof of \expref{Lemma}{lem:dense-subgraph} we shall use the
following simple claim, which is analogous to~\cite[Claim 5.3]{HK11}:

\begin{claim}
  \label{claim:small_support_means_small_mass}
  The following holds w.\,h.\,p.\ over $\G$: For any set $S\subset[n]$ of size $|S| < c_2 \ln n$,
  in any $\eps$-equilibrium $(x,y)$ of the above
  game satisfying $\|x_{[n]}\| \ge 1-t$ and $\|y_{[n]}\| \ge 1-t$, the
  probability mass the column (or row) player places on such a set
  $S$ is less than $\alpha + \eps + t$.
\end{claim}

\begin{proof}
  Fix one such set $S$. Then with high
  probability, there is a row $i^*$ of $B$ such that $B_{i^*,j} = 1$ for all
  $j \in S$.  Indeed, the probability that such a row does not exist is
  exactly
  \begin{equation}\label{eq:small_support_small_mass_claim}
  \left(1 - \beta^{-|S|}\right)^N < \exp(-n^{-c_2 \ln 1/\beta} N) =
  \exp(-n^{\ln 1/\beta})\,.
  \end{equation} In other words, if the column player places a probability mass
  $p$ on $S$, the row player can achieve payoff $p$ by deviating to
  $i^*$, whereas the payoff of the row player in $(x,y)$ is at most
  $\alpha (1-t) + t \le \alpha + t$.  Thus if $(x,y)$ is an
  $\eps$-equilibrium, $p$ is at most $\alpha + t + \eps$. Finally, to show that this holds for all such sets $S$, it is enough to take a union bound in equation~\eqref{eq:small_support_small_mass_claim}, so that (for sufficiently large $n$) the probability that the claim does not hold for at least one set $S$ is at most
  $$n^{c_2\ln n}\left(1 - \beta^{-|S|}\right)^N < 
  \exp(c_2(\ln n)^2-n^{\ln 1/\beta}) < \exp(-n^{\ln 1/\beta}/2)\,.$$
\end{proof}



%
%
%
%
%
%
%
%
%
%

\begin{proof}[Proof of \expref{Lemma}{lem:dense-subgraph}]
  Let $\tilde{x}, \tilde{y} \in [0,1]^n$ be the strategies $x$ and $y$
  conditioned on playing in $[n]$.  That is,
  $\tilde{x}=x_{[n]}/\|x_{[n]}\|$, and
  $\tilde{y}=y_{[n]}/\|y_{[n]}\|$.  The second given property of the
  strategy pair $(x,y)$ can be rephrased as $\tilde{x}^{\top} A
  \tilde{y} \ge 1-s$.
  
  Let
  $$S'_1 = \left\{ i \in [n]: e_i^{\top} A \tilde{y} \ge
  1-2\sqrt{s}/3\right\}.$$ To obtain a lower bound on the cardinality $|S'_1|$, we shall lower bound
  $\|x_{S'_1}\|$ by $\eps + \alpha$ and then appeal to the claim.  By the first given property of $(x,y)$, we have 
  $$\|x_{S'_1}\| \ge 1-t - \|x_{[n] \setminus S'_1}\| \ge 1 - t - \|\tilde{x}_{[n]\setminus  S'_1}\|\,.$$
  We can bound the last term, $\|\tilde{x}_{[n] \setminus S'_1}\|$, from above using Markov's inequality, viz.,
  $$
  \|\tilde{x}_{[n] \setminus S'_1}\| = \Pr_{i \sim \tilde{x}}\left[1 - e_i^{\top} A \tilde{y} > 2\sqrt{s} /3\right] \le \frac{s}{2 \sqrt{s} / 3} = 3\sqrt{s}/2\,.
  $$ Thus we have $\|x_{S'_1}\| \ge 1 - t - 3\sqrt{s}/2 \ge \eps +
  \alpha + t$ and so by \expref{Claim}{claim:small_support_means_small_mass}, $|S'_1| \ge c_2 \ln n$.  Truncate $S'_1$
  by taking any arbitrary subset $S_1\subseteq S'_1$ of cardinality
  $|S_1|=c_2 \ln n$.
  
  Now let $\overline{x} \in [0,1]^n$ be the uniform distribution over
  $S_1$.  Note that $\overline{x}^{\top} A \tilde{y} \ge
  1-2\sqrt{s}/3$.  Let
  $$
  S'_2 = \left\{j \in [n]: \overline{x}^{\top} A e_j \ge 5/9\right\}.
  $$ The argument to lower bound $|S'_2|$ is similar to the argument
  for $|S'_1|$.  We get
  $$
  \|\tilde{y}_{[n] \setminus S'_2}\| = \Pr_{j \sim \tilde{y}}\left[1 - \overline{x}^{\top} A e_j > 4/9\right] \le \frac{2 \sqrt{s}/3}{4/9} = 3\sqrt{s}/2
  $$
  and therefore
  $$\|y_{S'_2}\| \ge 1 - t - 3 \sqrt{s}/2 \ge \eps + \alpha + t\,,$$
  as desired.  Again, we can truncate $S'_2$ by taking a subset $S_2\subseteq S'_2$ of cardinality $|S'_2|=c_2\ln n$. By construction $d(S_1, S_2) \ge 5/9$, and we are done.
%
\end{proof}


\subsection{\texorpdfstring{Hardness for $\eps$ close to $1/2$}{Hardness for epsilon close to 1/2}}\label{sec:eps-hardness}

To obtain Theorem~1.2 the main requirement is to
set $\alpha = {1}/{2} + O(\eta)$.  The values of $\beta$ and $\gamma$ are
essentially irrelevant in this case---the only thing needed is
that $\beta$ is bounded away from both $0$ and $\alpha$ and that $\gamma
\le {1}/{2}$.

\begin{lemma}
  \label{lem:weight-in-A}
  For any $t \le 1/2$, let $\alpha = {1}/{2} + t$, $\gamma \le {1}/{2}$ and $\G$ be the game of
  \expref{Reduction}{red:reduction}.  Then for any pair of strategies
  $(x,y)$ with value at least $v_\G(x,y) \ge \alpha - t^2$ it holds
  that $\|x_{[n]}\|$ and $\|y_{[n]}\|$ are both at least $1 - t$.
\end{lemma}

\begin{proof}
  Let $p = \|x_{[n]}\|$ and $q = \|y_{[n]}\|$.  As the value of any
  outcome outside the $\alpha A$ block is at most ${1}/{2}$, we have that
  the value of $(x,y)$ is at most
  $$pq \alpha + (1-pq) \frac{1}{2} = t pq + \frac{1}{2}\,,$$
  so that if the value is at least $\alpha - t^2$ we have
  \[
  t pq + \frac{1}{2} \ge \alpha - t^2 \qquad\text{and}\qquad
  pq \ge \frac{\alpha - t^2 - 1/2}{t} = 1-t\,.
  \]
  Since $p, q \in [0,1]$, it follows that they are both at least
  $1-t$.
%
\end{proof}

\begin{observation}
  \label{obs:value-after-condition}
  Let $(x,y)$ be any pair of strategies with value $v_{\G}(x,y) \ge
  {1}/{2}$ and $\|x_{[n]}\| > 0$, $\|y_{[n]}\| > 0$.  Then
  $v_{\G|[n]}(x,y) \ge v_{\G}(x,y)$, provided that $\gamma \le {1}/{2}$.
\end{observation}

Plugging this into \expref{Lemma}{lem:dense-subgraph}, we can now easily
complete the proof of hardness for $\eps$ close to ${1}/{2}$.

\begin{theorem}[Detailed Statement of Theorem~1.2]
  \label{thm:eps-hardness-formal}
  For some universal constant $C$ and every $\eta>0$ there exist $\delta = \Omega(\eta^2)$, $\alpha
  \ge {1}/{2}$ such
  that the following holds.  Given a graph $G = (V,E)$ we can
  in randomized polynomial time construct a bimatrix game $\G$ with maximum average payoff $\alpha$ (over all strategy pairs) such
  that, if $G\sim\HC{n}{C\ln n}$, %
 the following holds (w.\,h.\,p.\ over $G$ and $\G$):
  \begin{description}
  \item[Completeness] %
%
    $\G$ has a Nash equilibrium
    $(x,y)$ with value $\alpha$.
  \item[Soundness]  %
%
 Given any $(\frac{1}{2}-\eta)$-equilibrium with value $\geq\alpha-\delta$, we can recover the hidden clique in polynomial time.
  \end{description}
\end{theorem}

\begin{proof}
  Given a graph $G = (V,E)$, we apply Reduction~\ref{red:reduction} with parameters as follows.
  For some $t>0$ to be determined momentarily, let $\alpha = {1}/{2} + t$,
  $\beta = \gamma = 1/3$, $\delta = t^2$.

  For the completeness, we shall show that having both players play
  uniformly over the hidden clique is an equilibrium.  For this to
  hold, we have to make sure that there is no row in $B$ with average
  value at least $\alpha$ in the positions corresponding to the
  clique.  By \expref{Lemma}{lemma:chernoff} and a union bound over all
  rows of $B$ we can bound the probability of this happening by
  $$N e^{-2 (\alpha-\beta)^2 C \ln n} \le N \cdot n^{-C/18}\,.$$ If $C$
  is a sufficiently large universal constant (\eg, $C = 18 \cdot
  (c_2+1) \ln 1/\beta + 1$) this probability is $o(1)$ and the
  completeness property follows.

  For the soundness, consider any $({1}/{2}-\eta)$-approximate equilibrium $(x,
  y)$ with value at least $v_{\G}(x,y) \ge \alpha - \delta=\alpha-t^2$.  By
  \expref{Lemma}{lem:weight-in-A} both $\|x_{[n]}\|$ and $\|y_{[n]}\|$ are
  at least $1-t$.  Furthermore, by
  Observation~\ref{obs:value-after-condition} we have $v_{\G|[n]}(x,y)
  \ge \alpha - \delta \ge \alpha(1 - 2t^2)$.

  Thus if $1 - 2t - 3t/\sqrt{2} \ge \alpha + \frac{1}{2} - \eta$ we can apply
  \expref{Lemma}{lem:dense-subgraph} and extract a dense bipartite
  subgraph of $G$ which can be plugged in to
  \expref{Lemma}{lemma:reconstruct} to obtain the hidden clique.  Setting
  $t = \eta/6$, the desired inequality holds and we get the result.
%
\end{proof}

